\section{Oral English}

\begin{table}[h]
\centering
\caption{Common English Oral Mistakes \& Corrections}
\label{tab:oral_mistakes}
\begin{tabular}{|c|p{4cm}|p{4cm}|c|p{4cm}|}
\hline
\textbf{No.} & \textbf{Wrong Expression ❌} & \textbf{Right Expression ✅} & \textbf{Error Type} & \textbf{Deep Analysis \& Memory Anchor} \\ \hline
1 & I very like this movie. & I am really into this movie. / I highly recommend it. & Collocation & \textbf{CN Trap}: "Very" cannot directly modify verbs like "like". \newline \textbf{Upgrade}: Use "be obsessed with" (addicted) or "it appeals to me" (attractive). \\ \hline
2 & Because AI is convenient, so we use it. & Because AI is convenient, we use it. / AI is convenient, so we use it. & Conjunction Logic & \textbf{Grammar Law}: "Because" and "So" cannot coexist in one sentence. \newline \textbf{Bonus}: "Due to the convenience of AI" for formal writing. \\ \hline
3 & The price of this book is expensive. & The book is expensive. / The price is high. & Subject-Verb Logic & \textbf{Physics}: Prices have "High/Low"; Items have "Expensive/Cheap". \newline \textbf{Native}: "It costs a fortune." \\ \hline
4 & With the development of society, ... & As society evolves, ... / With rapid technological advancement, ... & Template Phrasing & \textbf{Avoid}: This phrase sounds childish in PET/FCE exams. \newline \textbf{Tip}: Cut to the chase or use "As... unfolds". \\ \hline
5 & I think this is a good thing. & From my perspective, this brings more pros than cons. / It’s a \textbf{double-edged sword}. & Opinion Stating & \textbf{Upgrade}: Avoid simple "Good/Bad". \newline \textbf{Logic}: Use "Admittedly", "However", "Ultimately" to build a logical chain. \\ \hline
\end{tabular}
\end{table}
